Write $s_{k} = \sqrt{b_{k}}$, $S_{1} = \sum_{k}s_{k}$, $S_{2} = \sum_{k}s_{k}^{2}$, so $Q = S_{2}/S_{1}^{2}$ (Proposition \claimref{families.linear-conjecture-refuted-data}). Two facts of Theorem \claimref{synthesis.mass-ordered-dressing-cannot-spoil-koide} are used: (a) $Q$ is homogeneous of degree zero, its part (i); (b) its part (ii): $\partial Q/\partial s_{j} = 2(S_{1}s_{j} - S_{2})/S_{1}^{3}$, which is negative whenever $3s_{j} \le s_{\max}$, so raising a root that stays at most a third of the largest root strictly lowers $Q$.

\emph{(i)} By (a) take the largest root $s_{3} = 1$ and order $s_{1} \le s_{2} \le 1$. A direct computation gives
\[
Q(t, t, 1) - \tfrac23 = \frac{3(2t^{2} + 1) - 2(2t + 1)^{2}}{3(2t + 1)^{2}} = -\frac{2t^{2} + 8t - 1}{3(2t + 1)^{2}},
\]
positive for $0 \le t < t_{*}$ and zero at $t = t_{*}$, where $t_{*} = -2 + \sqrt{72}/4 = 3/\sqrt{2} - 2 \in (0.1213, 0.1214)$, below $\tfrac13$. Suppose $s_{2} < t_{*}$. Raising $s_{1}$ to $s_{2}$ keeps it below $\tfrac13 = s_{3}/3$, so by (b) $Q(s_{1}, s_{2}, 1) \ge Q(s_{2}, s_{2}, 1) > \tfrac23$. So $Q = \tfrac23$ forces $s_{2} \ge t_{*}$, that is $b_{2} \ge t_{*}^{2}b_{3}$ with $t_{*}^{2} = \tfrac92 - 6\sqrt{2} + 4 = \tfrac{17}{2} - 6\sqrt{2}$. If $s_{2} = t_{*}$ and $s_{1} < t_{*}$, the same step is strict, $Q > \tfrac23$; so equality needs $s_{1} = s_{2} = t_{*}$, and $(t_{*}^{2}, t_{*}^{2}, 1)$ has $Q = \tfrac23$. For a rational $r$, $r < t_{*}^{2}$ exactly when $r < \tfrac{17}{2}$ and $(\tfrac{17}{2} - r)^{2} > 72$, decided in rationals. The check confirms the identity on rational samples, the brackets $t_{*} \in (0.1213, 0.1214)$ and $t_{*}^{2} \in (0.01471, 0.01472)$, and $Q > \tfrac23$ on a rational grid of the square $s_{1}, s_{2} \le 0.1213$.

\emph{(ii)} From the stored quantities (GeV): $m_{c}/m_{t} = 1.2730/172.57$ and $0.619/171.7$, in the stated intervals, and the rational test of (i) shows both below $t_{*}^{2}$; $m_{s}/m_{b} = 0.0935/4.18$ and $0.055/2.89$ are above it. For $0 < x \le m_{c}$ the triple $(x, m_{c}, m_{t})$ has middle entry $m_{c}$ and largest $m_{t}$, so by (i) $Q \ne \tfrac23$ (in fact $Q > \tfrac23$; the check also certifies $Q(m_{c}, m_{c}, m_{t}) > \tfrac23$ with rational bounds on the roots). Let $b$ be a Koide bracket of the up sector with $D_{c} \ge D_{t}$. Then $b_{c}/b_{t} = (m_{c}/m_{t})(D_{t}/D_{c}) \le m_{c}/m_{t} < t_{*}^{2} < 1$. If $b_{u} \le b_{c}$, the triple $b$ has largest entry $b_{t}$ and middle entry $b_{c}$, against (i). So $b_{u} > b_{c}$. If the bracket order is kept, $b_{u} \le b_{c} \le b_{t}$, then (i) gives $b_{c} \ge t_{*}^{2}b_{t}$, that is $D_{t}/D_{c} \ge t_{*}^{2}m_{t}/m_{c}$; by the equality case of (i) this is the value $\rho$ at which $Q(\rho m_{c}, \rho m_{c}, m_{t}) = \tfrac23$, and the check certifies $Q > \tfrac23$ at $\rho = 1.99$, $Q < \tfrac23$ at $\rho = 2.00$ (own scale) and the same at $4.08$ and $4.09$ ($M_{Z}$), with $9\rho m_{c} \le m_{t}$ at the upper ends so that $Q(\rho m_{c}, \rho m_{c}, m_{t})$ is strictly decreasing in $\rho$ there by (b).

\emph{(iii)} By (a) normalise $D_{3} = 1$, so $b = (\rho_{1}m_{1}, \rho_{2}m_{2}, m_{3})$ with $\rho_{j} = D_{3}/D_{j}$; the class $D_{2} > D_{3} > D_{1}$ is $\rho_{2} < 1 < \rho_{1}$. Part (iv) of the stored theorem assumes only $\rho_{j} \in (0, R]$ for $j = 1, 2$ and $Rm_{j} \le m_{3}/9$; with $R = \rho_{1}$ these hold for $\rho_{1} \le R^{*}$, so its bound applies in the class. For the sharp bound, $f_{1}(\rho) = Q(\rho m_{1}, m_{2}, m_{3})$ is continuous and, by (b), strictly decreasing on $1 \le \rho \le m_{3}/(9m_{1})$; the check certifies $f_{1} > \tfrac23$ at the lower and $f_{1} < \tfrac23$ at the upper end of each stated interval, both inside that range, so $R_{1}$ is the unique crossing and lies in the interval. If $\rho_{1} < R_{1}$, then raising $\sqrt{\rho_{2}m_{2}}$ to $\sqrt{m_{2}}$ (allowed by (b), since $9m_{2} \le m_{3}$ and $9\rho_{1}m_{1} < m_{3}$) lowers $Q$, so $Q(b) \ge Q(\rho_{1}m_{1}, m_{2}, m_{3}) = f_{1}(\rho_{1}) > \tfrac23$, the last step by the strict decrease of $f_{1}$ (here $\rho_{1} < R_{1} \le m_{3}/(9m_{1})$); hence $\rho_{1} \ge R_{1}$. The class $D_{1} > D_{3} > D_{2}$, $\rho_{1} < 1 < \rho_{2}$, is the same argument with the roles of $m_{1}$ and $m_{2}$ exchanged and $f_{2}(\rho) = Q(m_{1}, \rho m_{2}, m_{3})$, strictly decreasing on $1 \le \rho \le m_{3}/(9m_{2})$, which contains the stated intervals. That both classes contain Koide brackets: the check certifies $Q < \tfrac23$ at $\rho_{2} = 0.99$ with $\rho_{1} = 9, 13, 2200, 5200$ (and $Q > \tfrac23$ at $\rho_{1} = 1$), so by continuity in $\rho_{1}$ there is a Koide bracket with $\rho_{2} = 0.99$ and $\rho_{1}$ between, in the class $D_{2} > D_{3} > D_{1}$; for $D_{1} > D_{3} > D_{2}$ the check certifies $Q < \tfrac23$ at $\rho_{1} = 0.99$ with $\rho_{2} = 2.4, 2.9, 10, 21$, and $Q > \tfrac23$ at $\rho_{2} = 1$, which gives a Koide bracket in that class in the same way. In the up sector a Koide bracket of the class has $b_{u}/b_{t} = \rho_{1}m_{u}/m_{t} \ge R_{1}m_{u}/m_{t}$, and the check certifies $R_{1}m_{u} > m_{c} > \rho_{2}m_{c}$ at both mass sets, so $b_{u} > b_{c}$, in agreement with (ii).

\emph{(iv)} If $D_{3} \le D_{j}$ for $j = 1, 2$, part (ii) of the stored theorem (with its ninth condition, true for the measured values) gives $Q(b) \ge Q(m) > \tfrac23$, by its part (iii); so $\max_{j}\rho_{j} > 1$, and by its part (iv) with $R = \max_{j}\rho_{j}$, $Q(b) \ge Q(Rm_{1}, Rm_{2}, m_{3})$, which exceeds $\tfrac23$ for $R < R^{*}$. If $\max_{j}\rho_{j} = R^{*}$ and $\rho_{1} \ne \rho_{2}$, one of them is below $R^{*}$ and by (b) the inequality is strict, $Q(b) > \tfrac23$; so the least value is reached only at $\rho_{1} = \rho_{2} = R^{*}$. The two remaining strict orderings are the classes of (iii), with their sharper bounds; the up-sector reversal is (ii). The intervals for $R^{*}$ are those of the stored theorem; the check recertifies their upper ends from the stored quantities ($Q < \tfrac23$ at $\rho_{1} = \rho_{2} = 1.89$, $2.20$, $8.68$, $17.56$) and confirms on a rational grid inside the class $D_{2} > D_{3} > D_{1}$ with $\rho_{1}$ below the lower ends that $Q > \tfrac23$.
